%-------------------------------------------------------------------------------
%	SECTION TITLE
%
% GENERATED FILE -- do not edit directly.
% Edit ../cv-data.yaml (research_cv_summary) and this file's template
% (../templates/research-cv-aboutme.tex.j2), then run
% `python3 generate_cvs.py` from the repo root.
%-------------------------------------------------------------------------------
\cvsection{Summary}


%-------------------------------------------------------------------------------
%	CONTENT
%-------------------------------------------------------------------------------
\begin{cvparagraph}

A quantum communications researcher focused on stochastic channel modeling, entanglement distribution, and fading analysis for microwave-based quantum systems essential for distributed quantum computing.
Studying continuous-variable quantum information approaches that bridge classical communication theory and quantum-statistical performance, with the goal of building infrastructure for tomorrow's quantum networks.
A former backend software engineer, bringing practical experience in distributed systems, API development, and scalable architecture to the design of robust quantum communication solutions.
Currently, a PhD researcher at \href{https://cqilab.khu.ac.kr/}{CQILAB, Kyung Hee University}, Republic of Korea, under the supervision of Professor Shin Hyundong.
Previously worked as a backend software engineer, delivering projects across billing systems, KYC, e-commerce, e-procurement, and HRIS domains using PHP, RabbitMQ, Redis, and MinIO.
\end{cvparagraph}

\newcommand{\lineintv}{\vspace{0.0cm}} % Not in use

\cventry
  {} % Degree
  {Highlights\vspace{-0.3cm}} % Institution
  {} % Location
  {} % Date(s)
  {
    \begin{cvitems} % Description(s) bullet points
      \item Published across top-tier venues in robotics (IEEE/RSJ IROS, IEEE ICRA, IEEE RA-L, RSS, IEEE T-RO, IJRR), autonomous driving (IEEE IV, IEEE T-IV), and artificial intelligence (CVPR, ICCV, AAAI, NeurIPS, IEEE T-PAMI).
      \item Always pursue \textit{clean, useful, reproducible} code to support the community:
        \begin{itemize}
          \item e.g., Maintainer of \href{https://github.com/MIT-SPARK/TEASER-plusplus}{TEASER++ \githubstars{2,100}} and \href{https://github.com/MIT-SPARK/Hydra}{Hydra \githubstars{773}} \& publisher of \href{https://github.com/url-kaist/patchwork-plusplus}{Patchwork++ \githubstars{842}} and \href{https://github.com/LimHyungTae/patchwork}{Patchwork \githubstars{555}}, \\ \href{https://github.com/MIT-SPARK/KISS-Matcher}{KISS-Matcher \githubstars{643}}, \href{https://github.com/MIT-SPARK/VGGT-SLAM}{VGGT-SLAM \githubstars{882}}, \href{https://github.com/LimHyungTae/ERASOR}{ERASOR \githubstars{494}}, and \href{https://github.com/MIT-SPARK/BUFFER-X}{BUFFER-X \githubstars{243}}
        \end{itemize}
      \item Academic honors \& awards:
        \begin{itemize}
          \item IEEE ICRA 2025 Outstanding Reviewer Award (Selected as one of only 5 awardees out of 7,641 reviewers)
          \item \href{https://sites.google.com/view/rsspioneers2024/participants}{RSS Pioneers 2024} (Out of 202 candidates, only 30 researchers were selected)
          \item 2022 IEEE RA-L Best Paper Award (among 1,100 papers, only 5 papers were selected)
        \end{itemize}
      \item Won multiple SLAM challenges and competitions while pursuing technologies for real-world applications:
        \begin{itemize}
          \item 1st prize again in the \href{https://construction-robots.github.io/\#challenge}{HILTI SLAM Challenge'24 in IEEE ICRA}
          \item 1st prize in \href{https://hilti-challenge.com/}{HILTI SLAM Challenge'23 in IEEE ICRA} among 63 international teams
          \item 2nd cash prize in \href{https://hilti-challenge.com/}{HILTI SLAM Challenge'22 in IEEE ICRA} (in total, 4th place)
          \item 1st prize in \href{https://hitachi-lg.com/}{Hitachi-LG Data Storage} LiDAR application competition, Republic of Korea
          \item Received \href{https://www.ces.tech/innovation-awards/honorees/2023/honorees/h/hi-bot-hologram-image-guide-robot.aspx}{CES'23 innovation award} through technology transfer on SLAM for mobile robots (collaborated with HILLS Robotics)
        \end{itemize}
      \item Professional experience:
        \begin{itemize}
          \item Postdoctoral associate of Laboratory for Information \& Decision Systems (LIDS), MIT (advisor: \href{https://lucacarlone.mit.edu/}{Prof. Luca Carlone})
          \item Visiting scholar of Univ. Bonn, Germany (advisor: \href{https://www.ipb.uni-bonn.de/people/cyrill-stachniss/index.html}{Prof. Cyrill Stachniss})
          \item In 2022, served as a SLAM part outside expert, CTO division of LG Electronics, Republic of Korea
          \item Research intern in vision/deep learning team at \href{https://www.naverlabs.com/}{NAVER LABS}, Republic of Korea
        \end{itemize}
    \end{cvitems}
  }


\cventry
  {My research interests include, but are not confined to:} % Degree
  {Field of Interest}%\vspace{-0.3cm}} % Institution
  {} % Location
  {} % Date(s)
  {
    \begin{cvitems} % Description(s) bullet points
      \item {Continuous-variable quantum information and quantum microwave communication}
      \item {Stochastic channel modeling for quantum wireless links}
      \item {Entanglement distribution and degradation over noisy channels}
      \item {Fading analysis in open-air and ground-to-ground microwave systems}
      \item {Rydberg atomic receivers for quantum signal detection}
      \item {Quantum computing and distributed quantum information systems}
      \item {Bridging classical communication theory with quantum-statistical system design}
      % \item {Analytical redundancy}
    \end{cvitems}
  }
